\documentclass[12pt]{article}
\usepackage[utf8]{inputenc}
\usepackage[T1]{fontenc}
\usepackage{hyperref}
\title{Lalo, Dongshanba Language Resources for K-12 Curriculum Development}
\author{Compiled by Claude Sonnet 5.5 (claude-sonnet-5-5) for Andr\'as Kornai}
\date{}
\begin{document}
\maketitle
\section{Introduction}
Dongshanba Lalo (ISO 639-3 yik) is a Loloish (Sino-Tibetan) language of western Yunnan, China, also called Western Yi. The Dongshan variety is spoken in eastern Weishan and in Dali, Midu, Yongping and Baoshan counties. Wikipedia reports about 320,000 speakers for the whole Lalo cluster. The Global Recordings Network gives a much higher figure, which looks inflated. Wikipedia classes Lalo as vulnerable, and a Xiaowan-area study found it definitely endangered. It is not known to be taught in schools.
\section{Primary Educational Resources}
No school primers, textbooks or programmes were found in the search.
\section{Linguistic and Reference Materials}
Overview of the Lalo language and its dialects: \url{https://en.wikipedia.org/wiki/Lalo_language}

\emph{A Grammar of Lalo} (Lund University) describes the West mountain subdialect of Longjie township, so it is not Dongshanba itself: \url{https://www.lunduniversity.lu.se/lup/publication/00a3c596-d8d1-4438-a7fa-abe79ceead0a}

Article on vitality and endangerment of Lalo: \url{https://benjamins.com/catalog/ltba.38.2.05tin}

No dedicated Dongshanba grammar or dictionary was found.
\section{Religious and Community Resources}
None were found in the search.
\section{Digital and Audio Resources}
The Global Recordings Network has a listing for yik, with alternative names including Dongshanba Lalo and Jiantou. The search did not show what recordings it holds: \url{https://globalrecordings.net/en/language/yik}
\section{Curriculum Development Context}
Little Dongshanba-specific material is available. Any curriculum would need to build on the Lalo description of a neighbouring variety and on fieldwork with local speakers.
\end{document}